\documentclass[../../main.tex]{subfiles}
\begin{document}

\begin{flushright}
\footnotesize\textit{【自動文字起こし・要確認】}
\end{flushright}

{\bf [解]}
$O$から底面へ下ろした垂線の足を$H$（正$n$角形$A_1\cdots A_n$の中心）とし，$A_1H=x$（$0<x<1$）とおく．対称性から，
\begin{align}
V_n=n\times(\text{三角錐 } O\text{-}HA_1A_2) \quad\cdots\text{①}
\end{align}
である．$\triangle OHA_1$は$H$が直角，斜辺$OA_1=1$だから，ピタゴラスの定理より
\begin{align}
OH=\sqrt{1-x^2} \quad\cdots\text{②}
\end{align}
また，$\angle A_1HA_2=2\pi/n$（$H$を中心とした隣接2頂点のなす角）で，二等辺三角形$A_1A_2H$（$HA_1=HA_2=x$）の面積は
\begin{align}
\triangle A_1A_2H=\dfrac{1}{2}x^2\sin\dfrac{2\pi}{n}=x^2\sin\dfrac{\pi}{n}\cos\dfrac{\pi}{n} \quad\cdots\text{③}
\end{align}
$OH$は底面に垂直だから，三角錐$O$-$HA_1A_2$の高さは$OH$であり，②③から
\begin{align}
(O\text{-}HA_1A_2)=\dfrac{1}{3}\sqrt{1-x^2}\cdot x^2\sin\dfrac{\pi}{n}\cos\dfrac{\pi}{n} \quad\cdots\text{④}
\end{align}

①④から
\begin{align}
V_n=\dfrac{n}{3}\sin\dfrac{\pi}{n}\cos\dfrac{\pi}{n}\cdot x^2\sqrt{1-x^2}
\end{align}

$x^2\sqrt{1-x^2}$を$x\in(0,1)$で最大化する．$t=x^2\in(0,1)$とおくと$x^2\sqrt{1-x^2}=\sqrt{t^2(1-t)}=\sqrt{t^2-t^3}$だから，$f(t)=t^2-t^3$の最大値を考えれば良い．$f'(t)=2t-3t^2=t(2-3t)$より，
\begin{center}
\begin{tabular}{c|ccccc}
$t$ & $0$ & & $2/3$ & & $1$ \\
\hline
$f'$ & & $+$ & $0$ & $-$ & \\
\hline
$f$ & & $\nearrow$ & & $\searrow$ &
\end{tabular}
\end{center}
より$t=2/3$（$x^2=2/3$）で最大．この時
\begin{align}
x^2\sqrt{1-x^2}=\sqrt{f(2/3)}=\dfrac{2}{3}\sqrt{1-\dfrac{2}{3}}=\dfrac{2}{3}\cdot\dfrac{1}{\sqrt{3}}=\dfrac{2\sqrt{3}}{9}
\end{align}
したがって，
\begin{align}
V_n=\dfrac{n}{3}\sin\dfrac{\pi}{n}\cos\dfrac{\pi}{n}\cdot\dfrac{2\sqrt{3}}{9}=\dfrac{2\sqrt{3}}{27}\,n\sin\dfrac{\pi}{n}\cos\dfrac{\pi}{n}
\end{align}

\textbf{(2)} $n\sin(\pi/n)=\pi\cdot\dfrac{\sin(\pi/n)}{\pi/n}\to\pi$（$n\to\infty$），$\cos(\pi/n)\to1$だから，
\begin{align}
\lim_{n\to\infty}V_n=\dfrac{2\sqrt{3}}{27}\cdot\pi\cdot1=\dfrac{2\sqrt{3}\,\pi}{27}
\end{align}


\newpage

\end{document}
